\subsection{概率经典问题}

\begin{结论}{半圆问题与 Wendel 定理}
圆周上 $n$ 个独立均匀点落在某个半圆内的概率 \重点{\frac{n}{2^{n-1}}}（以每个点为“最左端”的事件互斥，各概率 $2^{-(n-1)}$）。
推广（Wendel）：$\mathbb R^d$ 中 $n$ 个对称分布的随机向量位于某个过原点半空间的概率 $2^{-n+1}\sum_{k=0}^{d-1}\binom{n-1}k$；球面上（$d=3$）$\frac{n^2-n+2}{2^n}$。
\end{结论}

\begin{结论}{随机圆内接三角形是锐角三角形的概率}
三点均匀分布在圆周上：锐角三角形 $\iff$ 圆心在三角形内 $\iff$ 三点不在同一半圆内，概率 $1-\frac{3}{4}=\重点{\frac14}$。
四点时四面体（球面上）包含球心的概率 $\frac18$（Putnam 1992），用对径点配对的对称性论证：每点及其对径点 $2^3$ 种选择中恰有一种包含中心。
\end{结论}

\begin{结论}{Bertrand 悖论：“随机弦”有歧义}
单位圆中随机弦长于内接正三角形边长的概率：随机端点 $\frac13$，随机径向距离 $\frac12$，随机中点 $\frac14$。
\重点{\text{“随机”须指定测度}}；要求对平移、旋转、伸缩都不变的测度（Jaynes）挑出 $\frac12$（直线测度 $\dd p\,\dd\theta$），与 Buffon 实验一致。
\end{结论}

\begin{结论}{折棒成多边形}
单位棒随机折成 $n$ 段，能拼成多边形 $\iff$ 每段 $<\frac12$：
\[
\重点{P=1-\frac{n}{2^{n-1}}}
\]
（“某段 $\ge\frac12$”的 $n$ 个事件互斥，各概率 $2^{1-n}$，与半圆问题对偶）。$n=3$ 给三角形概率 $\frac14$。
\end{结论}

\begin{结论}{随机两点的平均距离}
单位区间 $\frac13$；单位正方形 $\frac{2+\sqrt2+5\ln(1+\sqrt2)}{15}\approx0.5214$；单位圆盘（半径 $1$）\重点{\frac{128}{45\pi}\approx0.9054}；单位立方体 $\approx0.6617$；单位球面上弦长 $\frac43$。
单位正方形内随机三角形平均面积 $\frac{11}{144}$；单位圆盘内 $\frac{35}{48\pi}$。距离分布用 Crofton 微分方程或卷积 $\int_K\mathbf 1_K(x)\mathbf 1_K(x+r)\dd x$。
\end{结论}

\begin{结论}{随机二次方程有实根的概率}
$x^2+Bx+C=0$，$B,C\sim U(0,1)$ 独立：$P(B^2\ge4C)=\int_0^1\frac{b^2}{4}\dd b=\重点{\frac1{12}}$。
$B,C\sim U(-1,1)$：$\frac{1}{2}+\frac{1}{24}=\frac{13}{24}$（$C<0$ 时必有实根）。$ax^2+bx+c$，$a,b,c\sim U(0,1)$ 独立：$P(b^2\ge4ac)=\frac{5+6\ln2}{36}\approx0.254$。这类问题统一化为区域的面积/体积。
\end{结论}

\begin{结论}{会面问题}
两人在 $[0,T]$ 内独立均匀到达，各等 $t$：相遇概率 $1-\bigl(1-\frac tT\bigr)^2$。等待时间不同 $t_1,t_2$ 时为 $1-\frac12\bigl(1-\frac{t_1}T\bigr)^2-\frac12\bigl(1-\frac{t_2}T\bigr)^2$。
\重点{\text{二维均匀分布的事件概率 }=\text{ 面积比}}；推广到 $n$ 人两两相遇等同高维多面体体积。
\end{结论}

\begin{结论}{Buffon--Laplace：网格上的投针}
长 $\ell$ 的针投到 $a\times b$ 矩形网格（$\ell<\min(a,b)$）上，与网格线相交的概率
\[
\重点{P=\frac{2\ell(a+b)-\ell^2}{\pi ab}}.
\]
平行线（$b\to\infty$）回到 $\frac{2\ell}{\pi a}$。期望交点数 $\frac{2\ell(a+b)}{\pi ab}$ 对弯曲的针也成立（Buffon 面条，期望的可加性）。
\end{结论}

\begin{结论}{有限次掷骰的最优停止}
最多掷 $n$ 次，可在任一次停下得到点数：$V_1=3.5$，$V_{k+1}=\E\max(X,V_k)$，
\[
\重点{V_2=4.25,\quad V_3=\frac{14}{3}\approx4.67},\quad V_4\approx4.94 .
\]
策略：剩 $k$ 次时，当前点数 $>V_{k-1}$ 才停。逆向归纳是有限期最优停止的一般方法（期权的二叉树定价同理）。
\end{结论}

\begin{结论}{St. Petersburg 悖论}
第 $k$ 次首次出正面得 $2^k$ 元：期望无穷，但没人愿出高价。对数效用：期望效用有限，等价价格约 $4$ 元（初始财富很小时）；庄家资金上限 $B$ 时期望约 $\log_2B+1$ 元（$B=10^9$ 时约 $31$）。
\重点{\text{期望值不是决策的全部}}：方差、效用、对手偿付能力都要考虑——Kelly 准则的思想源头。
\end{结论}

\begin{结论}{两个信封悖论}
“换信封的期望收益 $\frac12\cdot2x+\frac12\cdot\frac x2=1.25x$”的错误：把 $x$ 当作固定值，同时又认为另一信封是 $2x$ 与 $\frac x2$ 等可能——对任何正常先验这不可能对所有 $x$ 成立。
\重点{\text{只要金额先验分布可归一，换与不换的期望相同}}；先验期望无穷时（重尾）会出现“总想换”的病态。
\end{结论}

\begin{结论}{随机排列中两元素同圈的概率}
均匀随机排列中，元素 $1$ 所在圈长在 $1..n$ 上\重点文{均匀分布}；$1$ 与 $2$ 在同一圈的概率 $\frac12$；$k$ 个指定元素全在同一圈的概率 $\frac1k$。
最长圈长度 $L$：$P(L>\frac n2)\to\ln2$；$\frac{L}{n}$ 的极限分布为 Golomb--Dickman 型，$\E\frac Ln\to0.6243$。
\end{结论}

\begin{结论}{随机映射的结构}
$[n]\to[n]$ 的均匀随机映射：是排列的概率 $\frac{n!}{n^n}\approx\sqrt{2\pi n}\ee^{-n}$；循环点个数 $\sim\sqrt{\frac{\pi n}2}$；连通分支数 $\sim\frac12\ln n$；迭代轨道的尾长与圈长期望均 $\sim\sqrt{\frac{\pi n}8}$。
\重点{\text{Pollard }\rho\text{ 分解与哈希碰撞检测（Floyd 判圈）}}的复杂度 $\sim\sqrt n$ 由此而来。
\end{结论}

\begin{结论}{随机排列的峰与单调段}
$n$ 元随机排列中，内部位置 $i$（$2\le i\le n-1$）是局部极大的概率 $\frac13$，峰的期望数 \重点{\frac{n-2}{3}}。
最长递增段的期望 $\sim\frac{\ln n}{\ln\ln n}$（与最长递增子列 $2\sqrt n$ 截然不同）。“升降交替排列”的个数是 Euler zigzag 数，指数生成函数 $\sec x+\tan x$。
\end{结论}

\begin{结论}{Banach 火柴盒之外：首次领先与等待时间}
公平硬币，正面累计数首次超过反面的时刻 $T$：$P(T=2n+1)=\frac{C_n}{2^{2n+1}}$（$C_n$ 为 Catalan 数），$P(T<\infty)=1$ 但 \重点{\E T=\infty}。
有偏硬币（正面 $p>\frac12$）：$\E T=\frac{1}{2p-1}$。同样的生成函数 $\frac{1-\sqrt{1-4pqs^2}}{2qs}$ 给出所有首达分布。
\end{结论}

\begin{结论}{抽签公平性与“第 $k$ 个人”}
$n$ 张签中一张有奖，依次不放回抽取：每个人中奖概率都是 \重点{\frac1n}，与顺序无关（交换性）。
推广：袋中 $a$ 白 $b$ 黑依次取出，第 $k$ 次取到白球的概率恒为 $\frac{a}{a+b}$；最后一个球是白球的概率也是 $\frac a{a+b}$。用对称性而非逐步条件概率计算。
\end{结论}

\begin{结论}{连续成功与“热手”统计偏差}
$n$ 次公平投篮，统计“连中 $k$ 次之后的命中率”：即使每次独立，按场次平均的这一比例\重点{\text{偏低于 }\frac12}（Miller--Sanjurjo 2018）。
原因：条件在“之前有 $k$ 连中”上的选择偏差（短序列中连中之后的样本偏少且偏向失败）。重新分析后，经典“热手谬误”研究的数据反而支持热手效应存在。
\end{结论}

\参考{Mosteller, \emph{Fifty Challenging Problems in Probability}；Grimmett--Stirzaker, \emph{One Thousand Exercises in Probability}；Gardner 系列；Flajolet--Sedgewick 第 II 章；Feller I；MIT 18.600、18.S34。}
